% Gemeinsame Deklarationen: unveränderte Aussagen, Kontexte und IDs.
\newcommand{\CSBPartDef}{%
\FormulaDefDeltaK[Cantor--Bernstein-Teilmenge]{%
  \begin{aligned}[t]
    &F\colon A\inj B\dsep G\colon B\inj A\vdash{}\\[-2pt]
    &\CBPart{F}{G}\coloneqq
      \{x\in A\mid \forall X\in\powerset(A)\,\bigl(\\[-2pt]
    &\qquad ((A\setminus G[B])\subseteq X\land G[F[X]]\subseteq X)
      \rightarrow x\in X\bigr)\}
  \end{aligned}%
}{CantorBernsteinPartDef}{%
  \DeltaRow{Mengen}{A\dsep B\dsep X\dsep x}
  \DeltaRow{Funktionen}{F\dsep G}
  \DeltaRow{\textbf{Neues Symbol}}{}
  \DeltaRow{Cantor--Bernstein-Teilmenge}{\CBPart{F}{G}}
}
}

\newcommand{\CSBPartSubset}{%
\FormulaThmDeltaK[Die Cantor--Bernstein-Teilmenge liegt im Ausgangsbereich]{%
  F\colon A\inj B\dsep G\colon B\inj A
  \vdash \CBPart{F}{G}\subseteq A%
}{CantorBernsteinPartSubset}{%
  \DeltaRow{Mengen}{A\dsep B\dsep x}
  \DeltaRow{Funktionen}{F\dsep G}
}
}

\newcommand{\CSBPartContainsSeed}{%
\FormulaThmDeltaK[Der Startbereich liegt im Cantor--Bernstein-Teil]{%
  F\colon A\inj B\dsep G\colon B\inj A
  \vdash A\setminus G[B]\subseteq\CBPart{F}{G}%
}{CantorBernsteinPartContainsSeed}{%
  \DeltaRow{Mengen}{A\dsep B\dsep X\dsep x}
  \DeltaRow{Funktionen}{F\dsep G}
}
}

\newcommand{\CSBPartMinimal}{%
\FormulaThmDeltaK[Minimalität des Cantor--Bernstein-Teils]{%
  \begin{aligned}[t]
    &F\colon A\inj B\dsep G\colon B\inj A\dsep X\subseteq A\dsep
      A\setminus G[B]\subseteq X\dsep G[F[X]]\subseteq X\vdash{}\\[-2pt]
    &\CBPart{F}{G}\subseteq X
  \end{aligned}%
}{CantorBernsteinPartMinimal}{%
  \DeltaRow{Mengen}{A\dsep B\dsep X\dsep Y\dsep x\dsep z}
  \DeltaRow{Funktionen}{F\dsep G}
}
}

\newcommand{\CSBPartClosed}{%
\FormulaThmDeltaK[Abgeschlossenheit des Cantor--Bernstein-Teils]{%
  F\colon A\inj B\dsep G\colon B\inj A
  \vdash G[F[\CBPart{F}{G}]]\subseteq\CBPart{F}{G}%
}{CantorBernsteinPartClosed}{%
  \DeltaRow{Mengen}{A\dsep B\dsep X\dsep y}
  \DeltaRow{Funktionen}{F\dsep G}
}
}

\newcommand{\CSBPartFixedPoint}{%
\FormulaThmDeltaK[Fixpunktgleichung des Cantor--Bernstein-Teils]{%
  \begin{aligned}[t]
    F\colon A\inj B\dsep G\colon B\inj A\vdash
    \CBPart{F}{G}={}&(A\setminus G[B])\\[-2pt]
      &{}\cup G[F[\CBPart{F}{G}]]
  \end{aligned}%
}{CantorBernsteinPartFixedPoint}{%
  \DeltaRow{Mengen}{A\dsep B}
  \DeltaRow{Funktionen}{F\dsep G}
}
}

\newcommand{\CSBComplementIdentity}{%
\FormulaThmDeltaK[Komplementform der Fixpunktgleichung]{%
  F\colon A\inj B\dsep G\colon B\inj A
  \vdash G[B]\setminus G[F[\CBPart{F}{G}]]=A\setminus\CBPart{F}{G}%
}{CantorBernsteinComplementIdentity}{%
  \DeltaRow{Mengen}{A\dsep B\dsep y}
  \DeltaRow{Funktionen}{F\dsep G}
}
}

\newcommand{\CSBSecondBranchImage}{%
\FormulaThmDeltaK[Bild des zweiten Cantor--Bernstein-Zweigs]{%
  \begin{aligned}[t]
    F\colon A\inj B\dsep G\colon B\inj A\vdash
    G[B\setminus F[\CBPart{F}{G}]]
      =A\setminus\CBPart{F}{G}
  \end{aligned}%
}{CantorBernsteinSecondBranchImage}{%
  \DeltaRow{Mengen}{A\dsep B}
  \DeltaRow{Funktionen}{F\dsep G}
}
}

\newcommand{\CSBFixedInjections}{%
\FormulaThmDeltaK[Cantor--Schröder--Bernstein für feste Injektionen]{%
  F\colon A\inj B\dsep G\colon B\inj A
  \vdash \exists H\colon A\bij B%
}{CantorBernsteinFixedInjections}{%
  \DeltaRow{Mengen}{A\dsep B}
  \DeltaRow{Funktionen}{F\dsep G\dsep H}
}
}
